\section{lightmem architecture}
\label{sec:method}
Analogous to the human memory, we design LightMem as shown in Figure \ref{fig:architecture}, which consists of three light modules:
\textit{Light1} implements an efficient \textit{Sensory Memory Module} that selectively preserves salient information from raw input (\S\ref{pre_compress}),
\textit{Light2} realizes a topic-aware \textit{STM Module} for transient information processing (\S\ref{STM}),
and \textit{Light3} provides an \textit{LTM module} designed to minimize test time update latency (\S\ref{LTM}) with a sleep time update mechanism. 
\change{The overall pipeline framework of \textbf{LightMem}, its specific models, and comparisons with other memory frameworks are presented in Appendix~\ref{background_detail}.
The complexity analysis for LightMem's efficiency gains is in Section~\ref{sec:complexity_analysis}.}

\subsection{Light1: Cognitive-inspired sensory memory}
\label{pre_compress}
In long horizon interaction scenarios, such as user–assistant dialogues, a large portion of the information is redundant. 
Therefore, we design a \textit{Pre-Compressing Submodule} to eliminate redundant tokens, followed by the \textit{Topic Segmentation Submodule} that forms semantic topic-based segments for following faster and more accurate memory construction. 

% \noindent \textbf{Pre-Compressing Submodule.} This module leverages an off-the-shelf prompt compression model $p_{\phi}$ to remove reduntant tokens. The compression process can be modeled as a token-level binary classification problem: 

\noindent \textbf{Pre-Compressing Submodule.} This module leverages a compression model $\theta$ to eliminate redundant tokens, tailored for compatibility with the downstream memory construction phase: 
\[
\hat{\mathbf{x}} = \left\{ x_i \in \mathbf{x} \;\middle|\; P(\text{retain } x_i \mid \mathbf{x}; \theta) > \tau \right\}, 
\tau = \text{Percentile}\left( \{x_j\}, r \right),
\]
Following~\citet{xia2025tokenskip}, we use LLMLingua-2~\citep{pan2024llmlingua} as our compression model $\theta$. Let $\mathbf{x}$ be the raw input tokens, $\theta$ the model, and $r$ the compression ratio. The threshold $\tau$ is set to the $r$-th percentile of retention scores, keeping only tokens above $\tau$. 
For $P(\text{retain } x_i \mid \mathbf{x})$, we treat the compression process as a binary token classification task (``retain'' or ``discard''). For each token $x_i$ in a sequence $\mathbf{x}$, the model $\theta$ outputs a logit vector $\ell_i$, and the retention probability is given by:

\[
P(\text{retain } x_i \mid \mathbf{x}; \theta) = \mathrm{softmax}(\ell_i)_1,
\]

where the subscript $1$ denotes the ``retain'' class. Tokens with probabilities above a dynamic threshold are included in the compressed sequence.
% The compression process is then defined as:
% \[
% P(\text{retain } x_i \mid \mathbf{x}; \theta) = \mathrm{softmax}(\ell_i)_{\text{dim1}}, \text{ }\ell_i: x_i \text{ logits}, \text{ dim1 corresponds to retaining class.}
% \]
% The model is selected for three main reasons:
% (1) It is a compact BERT-based model, making it lightweight;
% (2) Its bidirectional attention captures semantics using full context;
% (3) It is trained on extensive task-agnostic data, including commonsense reasoning, math, and format-specific constraints, making it suitable for user-agent dialogue. 
In addition, \textbf{LightMem} can also employ more general generative LLM as the pre-compression model. 
We further implement a token filtering mechanism based on the cross-entropy between the model’s predicted distribution and the true token labels: 
\[
P(\text{retain } x_i \mid \mathbf{x}; \theta) = -\sum_{x_i \in \mathcal{V}} q(x_i) \log P(x_i \mid \mathbf{x}; \theta)
\]
where $q(x_i)$ denotes the true token label distribution.
Tokens with higher conditional entropy under a given context are more uncertain and less predictable, indicating greater informational uniqueness and a more critical role in semantic expression, such distinctive tokens are essential for subsequent memory construction and are therefore retained.

\noindent \textbf{Topic Segmentation Submodule.}
Existing works indicate that topic-granular input facilitates improved performance in memory systems~\citep{DBLP:conf/iclr/PanWJLCL0LZQ025, tan2025prospectretrospectreflectivememory}. 
As shown in Figure~\ref{fig:architecture}, \textbf{LightMem} maintains a sensory memory buffer to temporarily store information after pre-compression. 
When the accumulated information reaches the buffer's maximum capacity, a hybrid topic segmentation operation based on attention and similarity is triggered. 
We use the compression model $\theta$ and an embedding model to compute attention matrices and semantic similarities, respectively.
We define the final segmentation boundaries as the intersection of attention-based boundaries $\mathcal{B}_1$ and similarity-based boundaries $\mathcal{B}_2$: 
\[
\begin{aligned}
\mathcal{B}_1 &= \left\{ k \mid M_{k,k-1} > M_{k-1,k-2},\ M_{k,k-1} > M_{k+1,k},\ 1 < k < n \right\}, 
\end{aligned}
\]
\[
\mathcal{B}_2 = \Bigl\{\, k \;\big|\; \mathrm{sim}(s_{k-1}, s_k) < \tau ,1 \leq k < n\Bigr\}, \quad 
\mathcal{B} = \mathcal{B}_1 \cap \mathcal{B}_2.
\]
Specifically, dialogue scenarios possess natural semantic units, namely the conversational turn. 
We construct a turn-level attention matrix $M \in \mathbb{R}^{n \times n}$. 
$\mathcal{B}_1$ are identified as local maxima in the sequence $\{M_{k,k-1}\}$, i.e., the sub-diagonal elements of $M$ corresponding to attention between consecutive sentences. 
The detailed process of $\mathcal{B}_1$ and illustrative cases are provided in Appendix \ref{appen_topic_seg}. 
To mitigate attention sinks and dilution in attention-based methods, we compute semantic similarity between adjacent turns near each candidate boundary in $\mathcal{B}_1$. Boundaries with similarity below threshold $\tau$ form set $\mathcal{B}_2$, which helps determine the final topic boundaries $\mathcal{B}$.

\subsection{Light2: Topic-aware short-term memory}
\label{STM}
% 输入输出没写好
% 减少api调用次数，实现一个trade off
% 得到一个个topic segmentation之后，形成了一个{topic, 具体内容}的索引形式，我们在STM buffer中先基于topic的语义相近程度进行一个聚类，形成了一个{big topic, small topic, 具体内容}的三级索引结构，然后再调用大模型f summary得到这些topic以及具体内容的精简版，有了这些topic的指导，具体内容的summary/extract也会更准确。
% 如果是直接将多session直接输入确实可以有效减少后续api调用次数，但是会因为信息主题过多而导致后续summary得到的memory entry不准确，导致性能下滑，通过topic约束的输入粒度既将api调用次数减少到了极致，又维持了性能的稳定
After obtaining individual topic segments, forming an index structure of \{topic, message turns\}, where message turns = \{$user_i$, $model_i$\}. 
These are first placed into the STM buffer. When the token count in the buffer reaches a preset threshold, we invoke LLM $f_{\text{sum}}$ to generate concise summaries of every structure. 
The final index structure stored in LTM is \{topic, \{$sum_i$, $user_i$, $model_i$\}\}. 
\[
\mathrm{sum}_i = f_{\mathrm{sum}}\!\left(S_i\right), \quad 
S_i \subseteq \{\, \mathrm{user}_i, \mathrm{model}_i \,\}, \; S_i \neq \varnothing,
\]
\[
\mathrm{Entry}_i = \left\{\, \mathrm{topic},\; \mathbf{e}_i := \operatorname{embedding}(\mathrm{sum}_i),\; \mathrm{user}_i,\; \mathrm{model}_i \,\right\},
\]
\noindent where $\mathrm{Entry}_i$ denotes the memory entry to be stored in LTM.
Compared with inputting at the granularity of a single turn or session, directly feeding multiple sessions can reduce subsequent API calls but often introduces inaccurate memory entries due to excessive topic mixing, leading to performance degradation. 
In contrast, topic-constrained input granularity minimizes API calls to the greatest extent while preserving summarization accuracy and maintaining stable system performance. 

\subsection{Light3: Long-term Memory with Sleep-time Update}
\label{LTM}
% % 在线时软删除，离线更新的并行化
% and improve upon the previously required serial update procedure by enabling parallelized updates of memory entries. 

\noindent \textbf{Soft Updating at Test Time.} 
At test time, when memory entries arrive, LightMem directly inserts them into LTM with soft updates, thereby decoupling the update process from online inference. 
Due to real-time updates being converted to direct insertions, interaction latency is significantly reduced. 
After all entries are inserted or when an update trigger arrives, we compute an update queue for every entry in LTM. 
\[
\mathcal{Q}(e_i) = \operatorname{Top}_{k}\Big\{ (e_j, \mathrm{sim}(v_i, v_j)) \;\mid\; t_j \geq t_i, j \neq i \Big\}_{:n},
\]
where $e_i$ denotes the $i$-th memory entry with embedding $v_i$ and timestamp $t_i$, $\mathrm{sim}(\cdot,\cdot)$ is the similarity function, and $\operatorname{Top}_{k}\{\cdot\}_{:n}$ indicates selecting the top-$k$ most similar candidates, with the update queue ${Q}(e_i)$ length fixed at $n$. 
Consistent with existing work, we select the top-$k$ existing memory entries with the highest semantic similarity as potential update sources. 
On this basis, we further impose the constraint that only entries with later timestamps are allowed to update earlier ones ($t_j \geq t_i$), which is consistent with realistic temporal dynamics. 
Here, $\mathcal{Q}(e_i)$ denotes the queue of other entries that may update $e_i$. 
Since this process involves only similarity retrieval, it is fast and lightweight, and can be executed offline in parallel with online inference. 

\noindent \textbf{Offline Parallel Update.}
LightMem does not simply transfer online update latency to offline phases, it substantially reduces the overall update latency. 
The online update mechanism in existing memory frameworks enforces sequential updates, leading to a total latency that accumulates with each update. 
As shown in Figure~\ref{fig:architecture}, in LightMem, each memory entry maintains a global update queue, with each queue corresponding to a distinct $f_{\text{update}}$ operation. 
Since the update targets are independent across queues, updates can be executed in parallel, thereby greatly reducing the total latency. 